%\vspace{\sectionReduceTop}
\section{Related Work}
%\vspace{\sectionReduceBot}

\setlength{\epigraphwidth}{0.55\columnwidth}
\epigraph{Reality is something you rise above.}{\textit{Liza Minnelli}}

The availability of large-scale 3D scene datasets~\cite{armeni_cvpr16,Song2017,Chang2017} and community interest in active vision tasks led to a recent surge of work that resulted in the development of a variety of simulation platforms for indoor environments~\cite{Kolve2017,Brodeur2017,gupta_cvpr17,Savva2017,Wu2018,Anderson2018-Language,Xia2018,Yan2018,Puig2018}.
These platforms vary with respect to the 3D scene data they use, the embodied agent tasks they address, and the evaluation protocols they implement.

This surge of activity is both thrilling and alarming. On the one hand, it is clearly a sign of the interest in embodied AI across diverse research communities (computer vision, natural language processing, robotics, machine learning).
On the other hand, the existence of multiple differing simulation environments can cause fragmentation, replication of effort, and difficulty in reproduction and community-wide progress.
Moreover, existing simulators exhibit several shortcomings:
\begin{compactitem}[\hspace{1pt}--]
    \item Tight coupling of task (\eg navigation), simulation platform (\eg GibsonEnv), and 3D dataset (\eg Gibson). 
    Experiments with multiple tasks or datasets are impractical. 
    \item Hard-coded agent configuration (\eg size, action-space). Ablations of agent parameters and sensor types are not supported, making results hard to compare.
    \item Suboptimal rendering and simulation performance. Most existing indoor simulators operate at relatively low frame rates (10-100 fps), becoming a bottleneck in training agents and 
    making large-scale learning infeasible. Takeaway messages from such experiments 
    become unreliable -- has the learning converged to trust the comparisons? 
    \item Limited control of environment state. The structure of the 3D scene in terms of present objects cannot be programmatically modified (\eg to test the robustness of agents). 
\end{compactitem}

Most critically, work built on top of any of the existing platforms is hard to reproduce independently from the platform, and thus hard to evaluate against work based on a different platform, even in cases where the target tasks and datasets are the same.
This status quo is undesirable and motivates the Habitat effort.
We aim to learn from the successes of previous frameworks and develop a unifying platform that combines their desirable characteristics while addressing their limitations.
A common, unifying platform can significantly accelerate research by enabling code re-use and consistent experimental methodology.
Moreover, a common platform enables us to easily carry out experiments testing agents based on different paradigms (learned vs. classical) and generalization of agents between datasets.

The experiments we carry out contrasting learned and classical approaches to navigation are similar to the recent work of Mishkin et al.~\cite{mishkin2019benchmarking}.
However, the performance of the Habitat stack relative to MINOS~\cite{Savva2017} used in 
\cite{mishkin2019benchmarking} -- thousands vs.~one hundred frames per second -- 
allows us to evaluate agents that have been trained with significantly larger amounts of 
experience (75 million steps vs.~five million steps). 
The trends we observe demonstrate that learned agents can begin to match and outperform classical approaches when provided with large amounts of training experience.
Other recent work by Koijima and Deng~\cite{kojima2019learn} has also compared hand-engineered navigation agents against learned agents but their focus is on defining additional metrics to characterize the performance of agents and to establish measures of hardness for navigation episodes.
To our knowledge, our experiments are the first to train navigation agents provided with multi-month experience in realistic indoor environments and contrast them against classical methods.
